\chapter{Технологическая часть}

Данные раздел содержит запросы и их результаты. В каждой БЯМ создавался новый чат, для исключения влияния контекста. 
Каждая большая языковая модель получала на вход одинаковый стартовый запрос.

\textit{Передо мной поставлена следующая задача: Цель – разработать и реализовать программное
обеспечение на языке Python для извлечения данных из текстовых файлов, полученных из pdf-файлов с применением библиотеки PyPDF2,
с использованием регулярных выражений.
Для достижения цели выполнить следующие задачи.\\
1. Разработать регулярные выражения для решения задачи по варианту.\\
2. Реализовать функцию для поиска подстроки по варианту в pdf-файле с использованием разработанных регулярных выражений (привести
листинг как самой функции, так и списка необходимых для её функционирования библиотек, а также пример вызова функции), возвращающую кортеж вида: - значение, определяю-
щее наличие строки согласно заданному регулярному выражению ("истина/ложь"); - список
кортежей (если "ложь то список пустой) с двумя составляющими каждого кортежа —
строкой, найденной с помощью регулярного выражения, и координатами для определения
найденной строки в документе (минимально необходимо указать номер страницы и номер
строки документа, в которой найдено вхождение искомой строки. Непосредственно таблицы
и изображения строками не считать, но их названия считаются строками).
Также нужно реализовать функцию, которая находит все .pdf файлы в директории files и подает их на функции проверки. \\
3. Реализовать программное обеспечение, принимающее на вход путь к pdf-файлу, использующее функцию из предыдущего пункта.\\
Вариант: Проверка на наличие некириллических символов внутри кириллических токенов. \\
Я понял так, что нам необходимо среди строк файла найти кириллические токены
и определить вхождение в них некириллических символов.
Знаки препиная и числа входят в множество кириллических алфавитов.
Функция определяет координаты строк, в которых содержится такое вхождеие.
Путь к файлу будет передаваться через параметры функции. Комментарии не писать.}

\section{Реализация с помощью модель Qwen Coder}

При реализации с помощью модели Qwen-coder пришлось использовать 2 промта.
В первую итерацию, код выводил результат в JSON формате, а также допускал большое ложно-положительных срабатываний
из-за присутствия специсмволов, таких как: (), :, ; и т.д. Первая итерация представлена в листинге~\ref{list:qwen-1}.

Для корректировок был введен следующий промт \textit{Результат хороший, но нужно сделать красивый вывод результата.
Кроме этого, сочетание кириллического символа со спецсимволом это тоже, верный вариант. Учти это и перепиши код.}

После повторного запроса был получен код, решающий задачу. Он приведен в листинге~\ref{list:qwen-2}.

\section{Реализация с помощью модель DeepSeek}

После первой итерации DeepSeek сгенерировал результат, который представлен в листинге~\ref{list:deepseek-1}.
Результат работает, но облатает существенным недостатком, при нахождении рядом спецсимвола и числа программа выдавала ложно-положительный результат.
Был произведен следующий запрос для коррекции запроса:
\textit{Спецсимволы, которые идут рядом должны проходить проверку \\
Программа должна запускаться без аргументов и обладать одним сценарием, обработкой файлов в директории files, тоесть в main не нужно ветвление по параметрам запуска, в main() нужна только функция $process\_pdf\_files\_in\_directory$.}

После коррекции был предложин следующий вариант, который выполнял все необходимые функции. Вторая итерация представлена на листинге~\ref{list:deepseek-2}.

\section{Реализация с помощью модель ChatGPT 5.1}

ChatGPT 5.1 решил поставленную задачу за одну иттерацию.
Полученный код представлен в листинге~\ref{list:chatgpt-1}.

\section{Вывод}

В данном разделе были описаны запросы и результаты работы трёх больших языковых
моделей. Результат каждой модели был приведен в листингах.
Работа полученной программы протестирована на заранее заготовленных файлах. 
